\chapter{Abalone -- Rules of the Game}
The following rules have been copied from the original game manual \cite{AbaRules}.\index{Abalone!rules}
 
\subsection*{Object of the Game} To push your opponent's marbles off the hexagonal board.

\subsection*{How to win} The fist player to eject six of the opponent's marbles from the board wins the game.

\subsection*{Start of play}
Setup the black and white marbles as in Figure \ref{fig:startpos}.
% *** EPS-Graphic *** 
\begin{figure}[htbp!]
\begin{center}
\footnotesize
\caption{Starting position}
\vspace{0.5cm}

\includegraphics[width=0.4\columnwidth]{ruleseps/startpos.eps}
\label{fig:startpos}
\end{center}
\begin{center}
\footnotesize
This is the initial position for a 2 player game.
\end{center}
\end{figure}

\begin{itemize}
\item Choose who plays black and who will be white, black begins.
\item Players move in turn.
\item Once played, a move cannot be changed.
\end{itemize}

\subsection*{Moving}
Two different kind of moves can be distinguished, \emph{inline moves} and \emph{broadside moves}, see Figures \ref{fig:inline} and \ref{fig:broadside} respectively.
% *** EPS-Graphic *** 
\begin{figure}[htbp!]
\begin{center}
\footnotesize                            
\caption{{\it Inline} move}
\vspace{5mm}
\includegraphics[width=0.4\columnwidth]{ruleseps/inline.eps}
\label{fig:inline}
\end{center}
\begin{center}
\footnotesize
\begin{tabular}{l}
One, two, or three marbles are moved along\\
the direction in which they are aligned.
\end{tabular}
\end{center}
\end{figure}
% *** EPS-Graphic *** 
\begin{figure}[htbp!]
\begin{center}
\footnotesize                            
\caption{{\it Broadside} move}
\vspace{5mm}
\includegraphics[width=0.4\columnwidth]{ruleseps/broad.eps}
\label{fig:broadside}
\end{center}
\begin{center}
\footnotesize
Two or three marbles are moved ``parallel''.
\end{center}
\end{figure}
\begin{itemize}
\item At his turn a player may move in any of the six possible directions one, two, or three marbles, provided there is an adjacent free space.
\item When two or three marbles of the same colour are moved together, they must be moved in the same direction.
\item Not more than three marbles of the same colour may be moved in a single manoeuvre.
\item A single move may not be for more than one space at a time.
\end{itemize}

\subsection*{Sumito} 
\index{Abalone!Sumito}
A {\sc Sumito} is a numerical superiority allowing to push your opponent's marbles.
\begin{itemize}
\item In order to push your opponent's marbles, one of the three only {\sc Sumito} positions must be set up, see Figure \ref{fig:sumito}.
% *** EPS-Graphic *** 
\begin{figure}[htbp!]
\begin{center}
\footnotesize
\caption{The three only {\sc Sumito} positions}
\vspace{5mm}
\includegraphics[width=0.4\columnwidth]{ruleseps/sumito.eps}
\label{fig:sumito}
\end{center}
\begin{center}
\footnotesize
\begin{tabular}{l}
2 Push 1 {\sc Sumito} \\
3 Push 2 {\sc Sumito} \\
3 Push 1 {\sc Sumito}
\end{tabular}
\end{center}
\end{figure}

\item Opponent's marbles may only be pushed {\it inline}, when in contact and then only providing there is a free space behind the marble or group of two marbles in the defensive position. See Figure \ref{fig:impossible} for some examples of illegal pushes.

% *** EPS-Graphic *** 
\begin{figure}[htbp!]
\begin{center}
\footnotesize
\caption{Examples of impossible pushes}
\vspace{5mm}
\includegraphics[width=0.4\columnwidth]{ruleseps/impossib.eps}
\label{fig:impossible}
\end{center}
\begin{center}
\footnotesize
\begin{tabular}{l}
Black cannot push White because there is no free space behind the white group; \\
because there is an empty space between the black group and the white group;\\
because the marbles are not in a line.
\end{tabular}
\end{center}
\end{figure}

\item A player is never obliged to push.
\end{itemize}

\subsection*{Pac} 
\index{Abalone!Pac}
A {\sc Pac} is a position of balanced forces.
\begin{itemize}
\item The only three possible {\sc Pac} positions are the following.
  \begin{itemize}   
  \item 1 to 1 {\sc Pac}   
  \item 2 to 2 {\sc Pac}
  \item 3 to 3 {\sc Pac} 
  \end{itemize}
\item More than three marbles of the same colour and {\it in line} do not change the 3 to 3 score.
\item In a {\sc Pac} position neither side can push. The {\sc Pac} must be broken along a different line of attack, see Figure \ref{fig:pac}.
% *** EPS-Graphic *** 
\begin{figure}[htbp!]
\begin{center}
\footnotesize
\caption{How to break a {\sc Pac}}
\vspace{5mm}
\includegraphics[width=0.4\columnwidth]{ruleseps/pac.eps} 
\label{fig:pac}
\end{center}
\begin{center}
\footnotesize
The horizontal {\sc Pac} can only be broken by the two black marbles.
\end{center}
\end{figure}
\end{itemize}

\subsection*{Ejection}
A marble is ejected when it is pushed off the board, see Figure \ref{fig:eject}. The first player to eject six of the opponent's marbles wins the game.

% *** EPS-Graphic *** 
\begin{figure}[htbp!]
\begin{center}
\footnotesize
\caption{Eject a marble}
\vspace{5mm}
\includegraphics[width=0.4\columnwidth]{ruleseps/eject.eps}
\label{fig:eject}
\end{center}
\begin{center}
\footnotesize
\begin{tabular}{l}
Black can push a white marble off the board.
\end{tabular}
\end{center}
\end{figure}

\subsection*{Key Words}
\begin{description}
\item[Move]  Moving a group of marbles without pushing your opponent's.
\item[Sumito] Numerical superiority in line: {\it to set up a {\sc Sumito}} (Figure     \ref{fig:sumito}).
\item[Push] The action of pushing one or two of your opponent's marbles. This is only possible in a {\sc Sumito} position.
\item[Line] Two or three marbles that move along the direction in which they are aligned: {\it inline move}.
\item[Broadside] A sideways escape move by two or three marbles together: {\it broadside move}. \item[Pac] A position of numerical equality in line: {\it {\sc Pac} position}.
\item[Eject] Pushing a marble off the board, see Figure \ref{fig:eject}.
\item[Gambit] A voluntary sacrifice of a marble to achieve a later gain. \index{Abalone!Gambit}
\item[Covered] A position in which a marble is surrounded or the surrounded marble itself: {\it covered ball}, {\it covered position}.
\item[Det] A marble in a position from which it can be ejected and having no escape route: {\it {\sc Det} marble} (Figure \ref{fig:eject}). \index{Abalone!Det}
\end{description}

